\section{Certificates and their independent check}\label{sec:cert}

Every lower bound reported in Section~\ref{sec:exp} is the output of a
separate checker, not of the solver.  This section specifies what a
certificate contains, what the checker verifies, and what has to be trusted.

\paragraph{Contents.}  A certificate is a compressed NumPy archive with
\begin{itemize}
\item the rows with positive multiplier, in compressed sparse row form: for
row $i$ the pairs $p=\{u,v\}$ of its support, integer coefficients $a_{ip}$
and an integer right-hand side $b_i$, in the $x$-form $a_i^\top x\ge b_i$;
\item the multipliers $y_i\ge0$ as double-precision numbers;
\item the identity of the instance: the name and relative path of the raw file
as distributed by PACE or SNAP, its SHA-256, its format, the number of
vertices $n$ and edges $m$, and the SHA-256 of the canonical edge list, that
is, of $n$ followed by the lexicographically sorted pairs $(u,v)$ with $u<v$
as 64-bit integers.
\end{itemize}
The vertex numbering is fixed by the instance definition: a PACE file numbers
the vertices $1,\dots,n$; in a SNAP edge list the vertices are the identifiers
that occur in a kept edge, numbered in increasing order, a signed list keeps
only its positive rows, and self-loops and duplicate edges are dropped.  The
clustering whose cost is reported is archived next to the certificate, one
label per vertex in the same numbering.

\paragraph{The check.}  The checker (\texttt{experiments/check\_certificate.py},
\numCheckerLines{} lines of Python with Numba kernels) imports nothing from
the solver package.  Given the raw file and the certificate it
\begin{enumerate}
\item parses the raw file with its own reader and rejects the certificate
unless $n$, $m$, the file hash and the edge-list hash agree with the recorded
ones;
\item checks that every pair of every row is in $P$, i.e.\ is an edge or has a
common neighbour, by intersecting sorted adjacency lists;
\item checks that every row is valid for every clustering that separates the
far pairs.  A row on at most nine vertices is checked by enumerating all
partitions of its vertex set (at most $21{,}147$, as restricted growth
strings) that separate its far pairs and taking the minimum of the integer
left-hand side.  A larger row is accepted only if it is a star row: its
positive entries are the pairs $vt$, $t\in T$, of one centre $v$, and its
negative entries are pairs $R$ inside $T$.  Its minimum over far-separating
clusterings is at least $|T|-|R|-M$ with $M=1$ if $R$ contains every non-far
pair inside $T$ (then a far-free subset of $T$ is a clique of $R$) and $M$
the independence number of $(T,R)$ otherwise, computed exactly;
\item rounds each $y_i$ down to a multiple of $2^{-30}$ and evaluates
\begin{equation}\label{eq:checked}
\LB=\sum_i b_iy_i+\sum_{p\in P}\Bigl(\min\bigl\{0,\,c_p-\textstyle\sum_i a_{ip}y_i\bigr\}-\min\{0,c_p\}\Bigr)
\end{equation}
in integer arithmetic (64-bit accumulation under a checked overflow bound,
then arbitrary-precision integers), and returns $\lceil\LB\rceil$;
\item optionally evaluates the cost of the archived clustering on the parsed
instance, also in integers.
\end{enumerate}
Rounding down keeps $y\ge0$, so the rounded multipliers are again a dual
certificate; nothing else in~\eqref{eq:checked} is inexact.  The sum runs
over the pairs that occur in some row, since the other terms vanish.
Expression~\eqref{eq:checked} is~\eqref{eq:lb} with $K=|\Nt|=
-\sum_{p\in P}\min\{0,c_p\}$ absorbed.

\begin{proposition}\label{prop:checker}
If the checker accepts a certificate for an instance, then the value it
returns is at most $\OPT$ of that instance.
\end{proposition}
\begin{proof}
By steps 2 and 3, every row is supported on $P$ and satisfied by the cut
vector of every clustering that separates far pairs.  By
Lemma~\ref{lem:half} this includes every optimal clustering, and
Theorem~\ref{thm:anytime}(a) with the rounded multipliers gives $\LB\le\OPT$.
$\OPT$ is an integer.
\end{proof}

The inequality behind the proof, that~\eqref{eq:checked} is at most the cost
of every far-separating clustering for any rows valid for such clusterings,
is proved in Lean~4 (\texttt{far\_dual\_bound}, \texttt{far\_dual\_le\_opt};
Appendix~\ref{app:lean}), as is the fact that optimal clusterings separate far
pairs.  What remains to be trusted is the checker's reader, the enumeration
and the star formula of step~3, the integer evaluation, and the Python and
NumPy runtime.  The tests of the repository check the reader against the
solver's loader (same $n$, $m$ and edge hash on all instances we use), the
star formula against the enumeration on random star rows with at most nine
vertices, the integer cost of a clustering against the solver's, and the
rejection of certificates with a perturbed right-hand side, a negative
multiplier, a pair outside $P$, or a different instance.

\paragraph{Cost of checking.}  Checking is much cheaper than computing the
bound: the largest certificate we report, on \texttt{\numCheckMaxGraph}, has
\numCheckMaxRows{} rows and is checked in \numCheckMaxTime\,s
(Section~\ref{sec:exp-cert}).  One command re-checks every archived
certificate against freshly downloaded instance files and writes one line per
certificate.

\paragraph{What is not certified.}  The certificate proves a lower bound, and
the archived clustering an upper bound; together they bound the gap of the
reported solution.  The approximation guarantees of Theorems~\ref{thm:sparse}
and~\ref{thm:certiflip} are statements about the algorithms, not about a
particular run, and are not checked.  Running times are measured, not
certified.
